%!TeX program = lualatex
\documentclass[aspectratio=169, 10pt, t]{beamer}
\usepackage{luatexja}
\usepackage[no-math]{luatexja-fontspec}

% --- フォント設定 ---
\setmainjfont[BoldFont=BIZUDPGothic-Bold.ttf, YokoFeatures={JFM=prop}]{BIZUDPGothic-Regular.ttf}
\setsansjfont[BoldFont=BIZUDPGothic-Bold.ttf, YokoFeatures={JFM=prop}]{BIZUDPGothic-Regular.ttf}

\renewcommand{\familydefault}{\sfdefault}
\renewcommand{\kanjifamilydefault}{\gtdefault}

% --- パッケージ ---
\usepackage{graphicx}
\usepackage{tikz}
\usepackage{xcolor}
\usetikzlibrary{shadows, calc}

% --- 色彩設定 ---
\definecolor{TMUNavy}{RGB}{67, 102, 176}
\definecolor{AlertRed}{RGB}{204, 0, 0}
\definecolor{TextBlack}{RGB}{50, 50, 50}

% --- Beamer設定 ---
\setbeamercolor{structure}{fg=TMUNavy}
\setbeamercolor{title}{fg=TextBlack}
\setbeamercolor{frametitle}{fg=white, bg=TextBlack}
\setbeamercolor{normal text}{fg=TextBlack}
\setbeamercolor{itemize item}{fg=TextBlack}
\setbeamercolor{itemize subitem}{fg=TextBlack}

\setbeamertemplate{itemize item}{\scriptsize\raise1.2pt\hbox{\donotcoloroutermaths$\bullet$}}
\setbeamertemplate{itemize subitem}{\tiny\raise1.5pt\hbox{\donotcoloroutermaths$\blacktriangleright$}}

% ▼▼▼ 右下のボタンを消す設定 ▼▼▼
\setbeamertemplate{navigation symbols}{}

\setbeamersize{text margin left=0.5cm, text margin right=0.5cm}

\begin{document}

\begin{frame}[t]
    % --- ロゴとフッターの配置 (TikZ overlay) ---
    \begin{tikzpicture}[remember picture, overlay]
        % ロゴ (右上)
        \node[anchor=north east, xshift=0cm, yshift=0cm] at (current page.north east) {
            \includegraphics[height=1.0cm, keepaspectratio]{images/Asakalab_logo.png}
        };
        
        % ▼▼▼ フッター (右寄せに変更) ▼▼▼
        % anchor=south east (右下基準), current page.south east (ページの右下隅)
        % xshift=-0.5cm で右端から少し離しています
        \node[anchor=south east, xshift=-0.5cm, yshift=0.1cm] at (current page.south east) {
            \tiny \textcolor{TextBlack}{東京都立大学主催2025 年度電子情報システム工学科特別研究発表会，2026 年2 月16 日，東京都立大学日野キャンパス2 号館講義室A}
        };
    \end{tikzpicture}

    % --- ヘッダー ---
    \vspace{-0.7cm}
    \begin{center}
        {\large \textcolor{TextBlack}{ボ\hspace{0.1em}ト\hspace{0.01em}ルネックリンク帯域最大化を目的とした\\
        \vspace{0.05em}多制約最適経路問題における並列化改良型PSO}}
        
        \vspace{0.2em}
        {\footnotesize 22141073 渡瀬 遥己 \quad  指導教員： 朝香 卓也 }
        \vspace{-0.5cm} 
    \end{center}
    
    % 横線
    \par\noindent\rule{\linewidth}{0.4pt}
    \vspace{-0.5cm}

    % --- メインコンテンツ ---
    \small 
    \begin{columns}[T, onlytextwidth]
        
        % 左カラム
        \begin{column}{0.53\textwidth}
            \setlength{\leftmargini}{2.0em} 
            \vspace{0.2cm}
            \textcolor{TextBlack}{1.背景・課題}
            \vspace{-0.1cm}
            \begin{itemize} \setlength{\labelsep}{0.3em} \setlength{\itemsep}{1pt} \setlength{\parskip}{1pt}
                \item 背景： 5G/IoT普及に伴い、帯域・遅延等の複数条件を\\ \vspace{0.04cm}同時に満たす通信が必要
                \item 課題： 多制約最適経路問題は \textbf{\textcolor{AlertRed}{NP困難}} 
                    \begin{itemize} \setlength{\itemsep}{0pt}
                        \item 厳密解法は計算爆発、従来PSOは成功率が低い
                    \end{itemize}
            \end{itemize}

            \textcolor{TextBlack}{2.粒子群最適化 (PSO： Particle Swarm Optimization)}
            \vspace{-0.1cm}
            \begin{itemize} \setlength{\labelsep}{0.3em} \setlength{\itemsep}{1pt} \setlength{\parskip}{1pt}
                \item 鳥の群れがエサを探す行動を模倣した探索法
                \item 弱点： 群れ全体が局所解に捕まり、停滞
            \end{itemize}

            \textcolor{TextBlack}{3.提案：Restart PSO}
            \vspace{-0.1cm}
            \begin{itemize} \setlength{\labelsep}{0.3em} \setlength{\itemsep}{1pt} \setlength{\parskip}{1pt}
                \item リスタート戦略： 探索の停滞時、群れを\textbf{\textcolor{AlertRed}{強制的に再配置}}
                \item 効果： 局所解から強力に脱出し、新たな良い経路を発見
            \end{itemize}

            \textcolor{TextBlack}{4.評価結果}
            \vspace{-0.1cm}
            \begin{itemize} \setlength{\labelsep}{0.3em} \setlength{\itemsep}{1pt} \setlength{\parskip}{1pt}
                \item 成功率： 従来PSO 約10\% $\to$ 提案手法 \textbf{\textcolor{AlertRed}{約60\%}}
            \end{itemize}

        \end{column}

        % 右カラム
        \begin{column}{0.45\textwidth}
            \centering
            % 図1
            \includegraphics[width=1.0\linewidth, keepaspectratio]{images/MCOP1min.pdf}
            \vspace{0.1cm} 
            % 図2
            \includegraphics[width=1\linewidth, keepaspectratio]{images/fig3_4_restart_strategy_fixed_v66.png}
        \end{column}

    \end{columns}

\end{frame}
\end{document}